
\medskip
\noindent\textbf{5.4.}
\phantomsection
\label{sga3:VIIA:5.4}

The universal property of \(U_p(\mathfrak g)\) implies that
\(U_p(\mathfrak g)\) is functorial in \(\mathfrak g\). Every
homomorphism of restricted \(p\)-Lie algebras
\[
  \varphi:\mathfrak g\longrightarrow\mathfrak h
\]
induces a unique homomorphism of unital algebras \(U_p(\varphi)\)
such that
\[
  j_{\mathfrak h}\circ\varphi
  =
  U_p(\varphi)\circ j_{\mathfrak g}.
\]
Here are some examples.

\smallskip
\noindent\textit{a)}
If \(\mathfrak h=0\), then \(U_p(\mathfrak h)\) identifies with the
base ring, and \(U_p(\varphi)\) is an algebra homomorphism
\[
  \varepsilon_{\mathfrak g}:
  U_p(\mathfrak g)\longrightarrow R
\]
called the \emph{augmentation}.

\smallskip
\noindent\textit{b)}
Let now \(\mathfrak h=\mathfrak g^\circ\), the Lie algebra opposite
to \(\mathfrak g\). Thus \(\mathfrak g^\circ\) has the same
underlying module and the same symbolic \(p\)-th power as
\(\mathfrak g\), while its bracket is the negative of the bracket
of \(\mathfrak g\). We may identify
\(U_p(\mathfrak g^\circ)\) with the algebra opposite to
\(U_p(\mathfrak g)\). Moreover, the isomorphism
\[
  x\longmapsto-x
\]
from \(\mathfrak g\) onto \(\mathfrak g^\circ\) induces an
isomorphism
\[
  c_{\mathfrak g}:
  U_p(\mathfrak g)
  \longrightarrow
  U_p(\mathfrak g^\circ)
  \simeq
  U_p(\mathfrak g)^\circ.
\]
We call \(c_{\mathfrak g}\) the \emph{antipode} of
\(U_p(\mathfrak g)\).

\smallskip
\noindent\textit{c)}
Finally, let \(\mathfrak f\) and \(\mathfrak g\) be two restricted
\(p\)-Lie algebras, and let
\(\mathfrak h=\mathfrak f\times\mathfrak g\) be their product.
Its underlying \(R\)-module is the direct product, and
\[
  [(x,y),(x',y')]
  =
  ([x,x'],[y,y']),
  \qquad
  (x,y)^{(p)}
  =
  (x^{(p)},y^{(p)}).
\]
If
\[
  h_1:\mathfrak f\to\mathfrak k,
  \qquad
  h_2:\mathfrak g\to\mathfrak k
\]
are restricted \(p\)-Lie algebra homomorphisms satisfying
\[
  [h_1(x),h_2(y)]=0
  \quad
  (x\in\mathfrak f,\ y\in\mathfrak g),
\]
then
\[
  h_1+h_2:(x,y)\longmapsto h_1(x)+h_2(y)
\]
is a restricted \(p\)-Lie algebra homomorphism. Conversely, every
homomorphism \(\mathfrak f\times\mathfrak g\to\mathfrak k\) is of
this form. This characterizes \(\mathfrak f\times\mathfrak g\) by
a universal property.

For example, the maps
\[
  h_1:x\longmapsto j_{\mathfrak f}(x)\otimes1,
  \qquad
  h_2:y\longmapsto1\otimes j_{\mathfrak g}(y)
\]
induce a homomorphism \(h_1+h_2\) from
\(\mathfrak f\times\mathfrak g\) to the restricted \(p\)-Lie
algebra underlying
\[
  U_p(\mathfrak f)\otimes U_p(\mathfrak g).
\]
The universal properties of \(\mathfrak f\times\mathfrak g\) and
of the restricted enveloping algebras show that this homomorphism
extends to an isomorphism
\[
  \varphi:
  U_p(\mathfrak f\times\mathfrak g)
  \xrightarrow{\sim}
  U_p(\mathfrak f)\otimes U_p(\mathfrak g).
\]

\medskip
\noindent\textbf{Definition.}
\phantomsection
\label{sga3:VIIA:5.4-comultiplication-definition}
If \(\mathfrak f=\mathfrak g\), the diagonal map
\[
  \delta:x\longmapsto(x,x)
\]
from \(\mathfrak g\) to
\(\mathfrak g\times\mathfrak g\) induces a homomorphism from
\(U_p(\mathfrak g)\) to
\(U_p(\mathfrak g\times\mathfrak g)\). We denote by
\[
  \Delta_{\mathfrak g}:
  U_p(\mathfrak g)
  \longrightarrow
  U_p(\mathfrak g)\otimes U_p(\mathfrak g)
\]
the composite of this homomorphism with \(\varphi\).
\footnote{Thus
\(\Delta_{\mathfrak g}(x)=x\otimes1+1\otimes x\) for every
\(x\in\mathfrak g\); in particular,
\(\Delta_{\mathfrak g}\) is cocommutative.}

It is then easy to see that \(\Delta_{\mathfrak g}\), together
with the multiplication of \(U_p(\mathfrak g)\), makes
\(U_p(\mathfrak g)\) an \(R\)-coalgebra in groups (cf.\ \SGARefLink{sga3:VIIB:3.2:prop}{3.2}), with
augmentation \(\varepsilon_{\mathfrak g}\) and antipode
\(c_{\mathfrak g}\).

\medskip
\noindent\textbf{5.5.}
\phantomsection
\label{sga3:VIIA:5.5}
\footnote{The editors reorganized \S5.4.1 of the original into this
\S5.5. \SGARefLink{sga3:VIIA:5.5.1}{Proposition 5.5.1} combines results from \SGARefLink{sga3:VIIA:5}{Section~5} and
\SGARefLink{sga3:VIIA:3.2.3}{Proposition~3.2.3} and incorporates the original \SGARefLink{sga3:VIIA:7.3}{Lemma~7.3}; the
proof of Proposition 5.5.3(ii) expands the proof of
\((i)\Rightarrow(ii)\) in the original \SGARefLink{sga3:VIIA:7.4}{Theorem~7.4}.}

Let \(S\) be a scheme of characteristic \(p\). First let
\(\mathcal U\) be an \(\mathcal O_S\)-coalgebra in groups and put
\[
  G=\operatorname{Spec}^{*}\mathcal U.
\]
We saw in \SGARefLink{sga3:VIIA:3.2.3}{3.2.3} that, for every \(T\to S\),
\((\operatorname{Lie}G)(T)\) is the Lie subalgebra of
\(\Gamma(T,\mathcal U_T)\) consisting of primitive elements. If
\(x\) is primitive, then
\[
  \Delta(x^p)
  =
  x^p\otimes1+1\otimes x^p,
\]
since \(\binom pi\equiv0\pmod p\) for \(0<i<p\). Thus \(x^p\) is
again primitive. By \SGARefLink{sga3:VIIB:5.1:proposition}{5.1} and \SGARefLink{sga3:VIIB:5.2:theorem}{5.2}, the map \(x\mapsto x^p\) makes
\((\operatorname{Lie}G)(T)\) into an
\(\mathcal O(T)\)-restricted \(p\)-Lie algebra.

Now let \(\mathcal L\) be a quasi-coherent restricted
\(p\)-Lie algebra over \(\mathcal O_S\). As \(V\) ranges over the
open subsets of \(S\), the group-coalgebra structures defined
above on
\[
  U_p\bigl(\Gamma(V,\mathcal L)\bigr)
\]
induce on the associated sheaf
\(\mathcal U_p(\mathcal L)\) a structure of
\(\mathcal O_S\)-coalgebra in groups. Moreover, for every
\(S\)-scheme \(T\), there is an isomorphism
\[
  \mathcal U_p(\mathcal L_T)
  \xrightarrow{\sim}
  \mathcal U_p(\mathcal L)_T.
\]

Let \(\operatorname{Prim}\mathcal U_p(\mathcal L)\) denote the
subpresheaf which assigns to every open subset \(V\) the primitive
elements of \(\mathcal U_p(\mathcal L)(V)\); this is readily seen
to be a sheaf. As \(V\) ranges over the open subsets of \(S\), the
composites
\begin{center}
  \phantomsection
  \label{sga3:VIIA:5.5-primitive-sheaf-sequence}
  \(\displaystyle
\Gamma(V,\mathcal L)
  \xrightarrow{j}
\operatorname{Prim}\mathrm U_p\!\bigl(\Gamma(V,\mathcal L)\bigr)
  \longrightarrow
\operatorname{Prim}\mathcal U_p(\mathcal L)(V)
\)

\end{center}
define a morphism
\[
  \mathcal L
  \longrightarrow
  \operatorname{Prim}\mathcal U_p(\mathcal L),
\]
again denoted \(j\) or \(j_{\mathcal L}\). It in turn defines a
morphism of \(\mathcal O_S\)-restricted \(p\)-Lie algebras
\[
  \mathcal W(\mathcal L)
  \longrightarrow
  \operatorname{Prim}
    \mathcal W\bigl(\mathcal U_p(\mathcal L)\bigr)
\]
(cf.\ \SGARefLink{sga3:VIIA:3.2.3}{3.2.3}).

\medskip
\noindent\textbf{Proposition 5.5.1.}
\phantomsection
\label{sga3:VIIA:5.5.1}
\textit{Let \(\mathcal L\) be an
\(\mathcal O_S\)-restricted \(p\)-Lie algebra, locally free as an
\(\mathcal O_S\)-module. Then \(j_{\mathcal L}\) induces an
isomorphism of \(\mathcal O_S\)-restricted \(p\)-Lie algebras}
\[
  \mathcal W(\mathcal L)
  \xrightarrow{\sim}
  \operatorname{Prim}
    \mathcal W\bigl(\mathcal U_p(\mathcal L)\bigr).
\]

\begin{proof}
Let \(T\) be an \(S\)-scheme. In view of the identification
\[
  \mathcal U_p(\mathcal L_T)
  =
  \mathcal U_p(\mathcal L)_T,
\]
we must show that
\[
  \Gamma(T,\mathcal L_T)
  \longrightarrow
  \operatorname{Prim}
    \Gamma\bigl(T,\mathcal U_p(\mathcal L_T)\bigr)
\]
is bijective. Replacing \(S\) by \(T\), we reduce to \(T=S\); it
then suffices to prove that the morphism of sheaves
\[
  j_{\mathcal L}:
  \mathcal L
  \longrightarrow
  \operatorname{Prim}\mathcal U_p(\mathcal L)
\]
is an isomorphism. The question is local on \(S\), so suppose that
\(S=\operatorname{Spec}R\) and that \(\mathcal L\) is the sheaf
associated with an \(R\)-restricted \(p\)-Lie algebra \(L\) with a
totally ordered basis \((x_\alpha)\).

As in \SGARefLink{sga3:VIIA:5.3.3}{5.3.3}, let \(z_\alpha\) be the image of \(x_\alpha\) in
\(U=U_p(L)\). For a family
\(n=(n_\alpha)\), with \(0\leq n_\alpha\leq p-1\) and only
finitely many nonzero terms, put
\[
  z^{(n)}
  =
  \prod_\alpha\frac{z_\alpha^{n_\alpha}}{n_\alpha!},
\]
the product being taken in increasing order. Since
\[
  \Delta(z_\alpha)
  =
  z_\alpha\otimes1+1\otimes z_\alpha,
\]
one obtains
\[
  \Delta(z^{(n)})
  =
  \sum_r z^{(n-r)}\otimes z^{(r)},
\]
where the sum is over the finite set of \(r\) satisfying
\[
  0\leq r_\alpha\leq n_\alpha
  \quad\text{for every }\alpha.
\]
The \(z^{(n)}\), respectively the
\(z^{(n)}\otimes z^{(m)}\), form a basis of \(U\), respectively
of \(U\otimes U\). It follows that
\[
  \Delta(u)=u\otimes1+1\otimes u
\]
holds for \(u\in U\) if and only if \(u\) is a linear combination
of the \(z_\alpha\). This proves the proposition.
\end{proof}

\medskip
\noindent\textbf{Remark 5.5.2.}
\phantomsection
\label{sga3:VIIA:5.5.2}
Recall from \SGARefLink{sga3:VIIA:3.2.2}{3.2.2} and \SGARefLink{sga3:VIIA:3.2.3}{3.2.3} that the group \(S\)-functor
\[
  G=\operatorname{Spec}^{*}\mathcal U_p(\mathcal L)
\]
is very good and that
\[
  \operatorname{Lie}G
  =
  \operatorname{Prim}
    \mathcal W\bigl(\mathcal U_p(\mathcal L)\bigr).
\]
Thus \SGARefLink{sga3:VIIA:5.5.1}{Proposition 5.5.1} says that \(j_{\mathcal L}\) induces an
isomorphism
\[
  \mathcal W(\mathcal L)
  \xrightarrow{\sim}
  \operatorname{Lie}G.
\]

If, moreover, \(\mathcal L\) is locally free of finite rank, then
\(\mathcal U_p(\mathcal L)\) is finite locally free over
\(\mathcal O_S\) by \SGARefLink{sga3:VIIA:5.3.3}{5.3.3}. Hence
\(\operatorname{Spec}^{*}\mathcal U_p(\mathcal L)\) is represented
by the \(S\)-group
\[
  \mathfrak G_p(\mathcal L)
  =
  \operatorname{Spec}\mathcal U_p(\mathcal L)^{*}
\]
(cf.\ \SGARefLink{sga3:VIIA:3.2.2.1}{3.2.2.1}). We obtain the following more precise statement.

\medskip
\noindent\textbf{Proposition 5.5.3.}
\phantomsection
\label{sga3:VIIA:5.5.3}
\textit{Let \(\mathcal L\) be an
\(\mathcal O_S\)-restricted \(p\)-Lie algebra, locally free of
finite rank as an \(\mathcal O_S\)-module. Put}
\[
  \mathcal A=\mathcal U_p(\mathcal L)^{*},
  \qquad
  G=\mathfrak G_p(\mathcal L)=\operatorname{Spec}\mathcal A.
\]
\textit{Then:}

\smallskip
\noindent
\textit{\SGAInlineRefTarget{sga3:VIIA:5.5.3:item:i}(i) \(j_{\mathcal L}\) induces an isomorphism}
\[
  \mathcal W(\mathcal L)
  \xrightarrow{\sim}
  \operatorname{Lie}G
\]
\textit{of \(\mathcal O_S\)-restricted \(p\)-Lie algebras.}

\smallskip
\noindent
\textit{\SGAInlineRefTarget{sga3:VIIA:5.5.3:item:ii}(ii) Let \(\mathcal I\) be the augmentation ideal of
\(\mathcal A\) and put
\(\omega_G=\mathcal I/\mathcal I^2\) (cf.\ II, \SGARefLink{sga3:ii:remark:4-11-4}{4.11.4}). Then}
\[
  \omega_G
  \simeq
  \mathcal L^{*}
  =
  \operatorname{Hom}_{\mathcal O_S}
    (\mathcal L,\mathcal O_S).
\]
\textit{In particular, \(\omega_G\) is locally free of finite rank,
and \(\omega_{G/S}^{*}\simeq\mathcal L\).}

\begin{proof}
Part~\SGARefLink{sga3:vii:paragraph:5-1:equation:i}{(i)} follows from \SGARefLink{sga3:VIIA:5.5.2}{Remark~5.5.2}. We prove~\SGARefLink{sga3:vii:paragraph:5-1:equation:ii}{(ii)}. Let
\(\eta_{\mathcal U}\) and \(\varepsilon_{\mathcal U}\) denote the
unit and augmentation of
\(\mathcal U=\mathcal U_p(\mathcal L)\), and let
\(\eta_{\mathcal A}\) and \(\varepsilon_{\mathcal A}\) denote
those of \(\mathcal A\). Put
\[
  \mathcal J=\ker\varepsilon_{\mathcal U}.
\]
Then
\begin{equation}
  \tag{1}\label{sga3:vii:proposition:5-5-3:equation:1}
  \mathcal U
  =
  \eta_{\mathcal U}(\mathcal O_S)\oplus\mathcal J.
\end{equation}
Let \(\delta\) be the morphism defined by the diagram
\begin{center}
  \phantomsection
  \label{sga3:VIIA:5.5.3-augmentation-kernel-delta-square}
  \begin{tikzcd}[row sep=3.2em,column sep=4.2em]
\mathcal J \arrow[r,"\delta"] \arrow[d,swap,"\tau"]
  & \mathcal J\otimes\mathcal J \\
\mathcal U \arrow[r,"\Delta"]
  & \mathcal U\otimes\mathcal U \arrow[u,swap,"\pi"]
\end{tikzcd}

\end{center}
where \(\tau\) and \(\pi\) are the inclusion and projection
deduced from decomposition~\SGARefLink{sga3:vii:proposition:5-5-3:equation:1}{(1)}. There is an exact sequence
\begin{center}
  \phantomsection
  \label{sga3:VIIA:5.5.3-primitive-exact-sequence}
  \(\displaystyle
(\ast)\qquad
0 \longrightarrow \mathcal L^*
  \xrightarrow{j_{\mathcal L}}
\mathcal J
  \xrightarrow{\delta}
\mathcal J\otimes\mathcal J .
\)

\end{center}

By \SGARefLink{sga3:VIIA:5.3.3}{5.3.3}, the \(\mathcal O_S\)-module
\(\mathcal J/\mathcal L\) is locally free, and by \SGARefLink{sga3:VIIA:5.5.1}{5.5.1} the
preceding sequence remains exact after every base change.
Consequently, \(\delta\) induces an isomorphism from
\(\mathcal J/\mathcal L\) onto a locally direct summand \(Q\) of
\(\mathcal J\otimes\mathcal J\). By duality, the exact sequence
therefore becomes
\begin{center}
  \phantomsection
  \label{sga3:VIIA:5.5.3-dual-exact-sequence}
  \(\displaystyle
(\ast\ast)\qquad
0 \longleftarrow \mathcal L
  \xleftarrow{{}^{t}\!j_{\mathcal L}}
\mathcal I
  \xleftarrow{{}^{t}\!\delta}
\mathcal I\otimes\mathcal I .
\)

\end{center}

Decomposition~\SGARefLink{sga3:vii:proposition:5-5-3:equation:1}{(1)} corresponds by duality to
\begin{equation}
  \tag{2}\label{sga3:vii:proposition:5-5-3:equation:2}
  \mathcal A
  =
  \eta_{\mathcal A}(\mathcal O_S)\oplus\mathcal I.
\end{equation}
The transpose of \(\Delta\) is the multiplication
\[
  m_{\mathcal A}:
  \mathcal A\otimes\mathcal A
  \longrightarrow
  \mathcal A.
\]
Since \(\mathcal I\) is an ideal of \(\mathcal A\),
\(m_{\mathcal A}\) sends
\(\mathcal I\otimes\mathcal I\) into \(\mathcal I\). More
precisely, decomposition~\SGARefLink{sga3:vii:proposition:5-5-3:equation:2}{(2)} gives the commutative square
\begin{center}
  \phantomsection
  \label{sga3:VIIA:5.5.3-augmentation-ideal-multiplication-square}
  \begin{tikzcd}[row sep=3.1em,column sep=4.2em]
\mathcal I
  & \mathcal I\otimes\mathcal I \arrow[l,swap,"m'"] \arrow[d,"{}^{t}\!\pi"] \\
\mathcal A \arrow[u,"{}^{t}\!\tau"]
  & \mathcal A\otimes\mathcal A . \arrow[l,swap,"m_{\mathcal A}"]
\end{tikzcd}

\end{center}
This shows that the restriction \(m'\) of \(m_{\mathcal A}\) to
\(\mathcal I\otimes\mathcal I\) is the transpose of \(\delta\).
The dual exact sequence therefore yields
\[
  \mathcal I/\mathcal I^2
  \simeq
  \mathcal L^{*},
\]
which proves the proposition.
\end{proof}

